\section{Introduction}\label{sec:intro}

The major-element chemistry of the lunar surface records how the crust formed and evolved. Highland plagioclase is rich in Al and Ca; mare basalts are rich in Fe, Mg and Ti. Global maps of these elements constrain magma-ocean crystallisation, mare volcanism and impact mixing \citep{Heiken1991}. Orbital gamma-ray spectroscopy has provided global maps of Fe, Ti, Th and K, and model-dependent estimates of Mg, Al, Si and Ca, at a resolution of hundreds of kilometres \citep{Lawrence1998,Prettyman2006}.

X-ray fluorescence (XRF) excited by solar X-rays offers better spatial resolution and direct sensitivity to Mg, Al and Si \citep{Jenkins1999,Beckhoff2006}. The Chandrayaan-1 X-ray Spectrometer (C1XS) demonstrated the technique during a period of very low solar activity, when only a few flares gave usable spectra \citep{Grande2009,Narendranath2011}. Its successor on Chandrayaan-2, the Large Area Soft X-ray Spectrometer (CLASS), has a geometric area of \SI{64}{cm^2} and a footprint of $\approx$\SI{12.5}{km} \citep{Narendranath2014,Narendranath2020,Pillai2021}. It has operated through the rise of solar cycle 25, and its Level-1 spectra are public.

Turning XRF line intensities into composition requires the spectrum of the incident solar X-rays. The fluorescence yield of each element depends on the solar flux above its absorption edge. A softer or harder solar spectrum therefore changes line ratios such as Mg/Si or Ca/Si even when the surface composition is unchanged. The Chandrayaan-2 Solar X-ray Monitor (XSM) measures this spectrum \citep{Mithun2020}, and the instrument team's analyses combine CLASS spectra with simultaneous XSM data. For users without an XSM-based workflow, the usual fallback is to assume a reference solar spectrum or to rely on line ratios. Ratios do not remove the dependence on the solar spectrum, because each element samples a different part of it.

In this paper we develop and validate a CLASS processing chain that can be reproduced entirely from public data, without XSM spectra. We take the solar excitation from the GOES X-Ray Sensor: its two-channel temperature, applied to isothermal CHIANTI spectra, gives a time-resolved solar spectrum for every CLASS exposure. Our specific contributions are:
\begin{enumerate}
\item a simplified fundamental-parameter forward model (Sherman primary fluorescence with matrix attenuation and incidence geometry, scattered solar spectrum, binned Gaussian detector response), fitted to raw counts by Poisson maximum likelihood with per-spectrum uncertainties, goodness of fit and quality flags;
\item the identification of a solar-driven instrumental Al K$\alpha$ component in CLASS spectra, present only when the spacecraft is sunlit, and a two-template background model that separates it from the particle background;
\item a quantified error budget from synthetic recovery tests, including the bias caused by errors in the solar temperature, which shows which element ratios can and cannot be recovered without a measured solar spectrum;
\item global Mg/Si and Al/Si maps from \NMonths{} monthly archives, with observation-count and uncertainty layers, validated against returned-sample compositions at nine landing sites and against Lunar Prospector GRS maps, using analysis choices fixed before validation.
\end{enumerate}
We report element ratios relative to Si, not absolute abundances, and we state explicitly which quantities the data cannot constrain.
